\chapter{Sciatica}

\section*{Definition}
Pain radiating along the path of the sciatic nerve from lower back through hips and legs down to the foot.

\section*{Pathophysiology}
Caused by compression or irritation of the sciatic nerve typically at lumbar spine level; burning shooting quality distinguishes from mechanical back pain; may include numbness tingling in affected leg distribution.

\section*{Etiology}
\begin{itemize}
  \item Herniated lumbar disc compressing nerve root
  \item Spinal stenosis narrowing spinal canal
  \item Piriformis syndrome muscle entrapment
  \item Degenerative disc disease nerve compression
  \item Spondylolisthesis vertebral slippage
\end{itemize}

\section*{Indications for Evaluation}
Seek care if:
\begin{itemize}
  \item Severe or worsening symptoms
  \item Symptoms lasting more than sudden severe pain loss bladder bowel control saddle anesthesia requires immediate evaluation
  \item Accompanied by other concerning signs
\end{itemize}

\section*{Related Symptoms}
\begin{itemize}
  \item Lower back pain
  \item Numbness in leg
  \item Muscle weakness
  \item Radiating pain pattern
  \item Burning/shooting quality
\end{itemize}
\textbf{Note:} Always consider serious underlying conditions when evaluating persistent sciatica.


\section*{Self-Care / Home Management}
\begin{itemize}
  \item Apply the RICE protocol: Rest, Ice (15-20 min every 2-3 hours), Compression with elastic bandage, Elevation above heart level.
  \item Gradual return to activity with gentle stretching once acute pain subsides.
  \item Over-the-counter anti-inflammatory medications (ibuprofen, naproxen) can reduce pain and swelling.
  \item Avoid activities that worsen symptoms until healing is well underway.
  \item Seek evaluation for severe pain, inability to bear weight, joint deformity, or symptoms lasting beyond 2 weeks.
\end{itemize}

\section*{Diagnostic Approach}
\begin{itemize}
  \item History focuses on mechanism of injury, onset (acute vs insidious), location, aggravating/relieving factors, and prior episodes.
  \item Physical examination includes inspection (swelling, bruising, deformity), palpation, range-of-motion, and strength testing.
  \item Imaging (X-ray for fracture, ultrasound for tendon/ligament injury, MRI for soft tissue detail) guides diagnosis.
  \item Functional assessment evaluates ability to perform daily activities and sports-specific movements.
\end{itemize}

\section*{Treatment / Management Overview}
\begin{itemize}
  \item Acute injuries benefit from RICE, activity modification, and gradual rehabilitation exercises.
  \item Physical therapy is central to recovery, focusing on range-of-motion, strengthening, and proprioceptive training.
  \item Corticosteroid injections may be considered for significant inflammatory conditions (tendinitis, bursitis).
  \item Orthopedic referral is indicated for complete tears, fractures, joint instability, or failure of conservative therapy for 6-8 weeks.
\end{itemize}

\clearpage
